\begin{ex}
Điền số thích hợp vào ô trống: $4 + \dots = 9$.
\shortans{5}
\loigiai{
Muốn tìm số hạng chưa biết, ta lấy tổng trừ đi số hạng đã biết: $9 - 4 = 5$.
\textbf{Đáp số:} $5$.
}
\end{ex}

\begin{ex}
Điền số thích hợp vào ô trống: $\dots + 3 = 10$.
\shortans{7}
\loigiai{
Muốn tìm số hạng chưa biết, ta lấy: $10 - 3 = 7$.
\textbf{Đáp số:} $7$.
}
\end{ex}

\begin{ex}
Tìm số tự nhiên $a$ biết rằng: $a + 2 = 8$.
\shortans{6}
\loigiai{
Số cần tìm là: $a = 8 - 2 = 6$.
\textbf{Đáp số:} $6$.
}
\end{ex}

\begin{ex}
Tính giá trị của biểu thức: $3 + 4 + 2$.
\shortans{9}
\loigiai{
Thực hiện từ trái sang phải: $3 + 4 = 7$; sau đó $7 + 2 = 9$.
\textbf{Đáp số:} $9$.
}
\end{ex}

\begin{ex}
Tính giá trị của biểu thức: $1 + 5 + 4$.
\shortans{10}
\loigiai{
Thực hiện từ trái sang phải: $1 + 5 = 6$; sau đó $6 + 4 = 10$.
\textbf{Đáp số:} $10$.
}
\end{ex}

\begin{ex}
Lan có 3 bông hoa, Mai có 4 bông hoa, Cúc có 2 bông hoa. Hỏi cả ba bạn có tất cả bao nhiêu bông hoa?
\shortans{9}
\loigiai{
Tổng số hoa của cả ba bạn là: $3 + 4 + 2 = 9$ (bông hoa).
\textbf{Đáp số:} $9$.
}
\end{ex}

\begin{ex}
Trong bể có 5 con cá vàng. Nam thả thêm 2 con cá bảy màu và 3 con cá mây. Hỏi lúc này trong bể có bao nhiêu con cá?
\shortans{10}
\loigiai{
Số cá có trong bể lúc này là: $5 + 2 + 3 = 10$ (con cá).
\textbf{Đáp số:} $10$.
}
\end{ex}

\begin{ex}
Hùng có 4 viên bi. Dũng cho Hùng thêm 3 viên bi nữa. Sau đó bố mua thêm cho Hùng 2 viên bi. Hỏi Hùng có tất cả bao nhiêu viên bi?
\shortans{9}
\loigiai{
Số viên bi Hùng có tất cả là: $4 + 3 + 2 = 9$ (viên bi).
\textbf{Đáp số:} $9$.
}
\end{ex}

\begin{ex}
Một cây có 4 quả bưởi, cây bên cạnh có nhiều hơn 2 quả bưởi. Hỏi cả hai cây có tất cả bao nhiêu quả bưởi?
\shortans{10}
\loigiai{
Cây thứ hai có số quả bưởi là: $4 + 2 = 6$ (quả).\nCả hai cây có số quả bưởi là: $4 + 6 = 10$ (quả bưởi).
\textbf{Đáp số:} $10$.
}
\end{ex}

\begin{ex}
Tổ Một trồng được 3 cây, tổ Hai trồng được số cây bằng tổ Một, tổ Ba trồng được 4 cây. Hỏi cả ba tổ trồng được bao nhiêu cây?
\shortans{10}
\loigiai{
Tổ Hai cũng trồng được 3 cây.\nCả ba tổ trồng được số cây là: $3 + 3 + 4 = 10$ (cây).
\textbf{Đáp số:} $10$.
}
\end{ex}

\begin{ex}
Tìm số thích hợp điền vào dấu chấm: $2 + 5 = 1 + \dots$.
\shortans{6}
\loigiai{
Vế trái: $2 + 5 = 7$.\nĐể vế phải bằng 7 thì số cần điền là: $7 - 1 = 6$.
\textbf{Đáp số:} $6$.
}
\end{ex}

\begin{ex}
Tìm số thích hợp điền vào dấu chấm: $3 + \dots = 4 + 4$.
\shortans{5}
\loigiai{
Vế phải: $4 + 4 = 8$.\nĐể vế trái bằng 8 thì số cần điền là: $8 - 3 = 5$.
\textbf{Đáp số:} $5$.
}
\end{ex}

\begin{ex}
Có bao nhiêu cặp hai số tự nhiên có tổng bằng 9 (không kể thứ tự)?
\shortans{5}
\loigiai{
Các cặp số có tổng bằng 9 là: $(0, 9), (1, 8), (2, 7), (3, 6), (4, 5)$. Có tất cả $5$ cặp.
\textbf{Đáp số:} $5$.
}
\end{ex}

\begin{ex}
Có bao nhiêu cặp hai số tự nhiên khác 0 có tổng bằng 10 (không kể thứ tự)?
\shortans{5}
\loigiai{
Các cặp số khác 0 có tổng bằng 10 là: $(1, 9), (2, 8), (3, 7), (4, 6), (5, 5)$. Có tất cả $5$ cặp.
\textbf{Đáp số:} $5$.
}
\end{ex}

\begin{ex}
Số nào cộng với 0 thì có kết quả bằng 8?
\shortans{8}
\loigiai{
Bất kỳ số nào cộng với 0 cũng bằng chính nó. Số đó là số $8$.
\textbf{Đáp số:} $8$.
}
\end{ex}

\begin{ex}
Tìm một số biết rằng lấy số đó cộng với 5 thì được số lớn nhất có một chữ số.
\shortans{4}
\loigiai{
Số lớn nhất có một chữ số là 9.\nSố cần tìm là: $9 - 5 = 4$.
\textbf{Đáp số:} $4$.
}
\end{ex}

\begin{ex}
Tìm một số biết rằng lấy số đó cộng với 2 rồi cộng thêm 3 thì được số 10.
\shortans{5}
\loigiai{
Ta có: $2 + 3 = 5$.\nSố cần tìm là: $10 - 5 = 5$.
\textbf{Đáp số:} $5$.
}
\end{ex}

\begin{ex}
Hai số tự nhiên liên tiếp có tổng bằng 9. Hỏi số lớn hơn là số mấy?
\shortans{5}
\loigiai{
Hai số tự nhiên liên tiếp có tổng bằng 9 là 4 và 5 ($4 + 5 = 9$). Số lớn hơn là số $5$.
\textbf{Đáp số:} $5$.
}
\end{ex}

\begin{ex}
Hai số tự nhiên liên tiếp có tổng bằng 7. Hỏi số bé hơn là số mấy?
\shortans{3}
\loigiai{
Hai số tự nhiên liên tiếp có tổng bằng 7 là 3 và 4 ($3 + 4 = 7$). Số bé hơn là số $3$.
\textbf{Đáp số:} $3$.
}
\end{ex}

\begin{ex}
Một hàng gồm 4 bạn nam và 5 bạn nữ. Hỏi hàng đó có tất cả bao nhiêu bạn?
\shortans{9}
\loigiai{
Tổng số bạn trong hàng là: $4 + 5 = 9$ (bạn).
\textbf{Đáp số:} $9$.
}
\end{ex}

\begin{ex}
Bình có 2 cái bút, An có nhiều hơn Bình 3 cái bút. Hỏi cả hai bạn có bao nhiêu cái bút?
\shortans{7}
\loigiai{
Số bút của An là: $2 + 3 = 5$ (cái bút).\nCả hai bạn có tất cả số bút là: $2 + 5 = 7$ (cái bút).
\textbf{Đáp số:} $7$.
}
\end{ex}

\begin{ex}
Một xe buýt ban đầu có 4 hành khách. Đến trạm dừng có thêm 3 người lên xe, trạm tiếp theo lại có thêm 2 người lên xe. Hỏi trên xe có tất cả bao nhiêu hành khách?
\shortans{9}
\loigiai{
Tổng số hành khách trên xe là: $4 + 3 + 2 = 9$ (người).
\textbf{Đáp số:} $9$.
}
\end{ex}

\begin{ex}
Tìm số tự nhiên $x$ sao cho $x + 3 < 5$. Có bao nhiêu giá trị của $x$?
\shortans{2}
\loigiai{
Vì $x$ là số tự nhiên và $x + 3 < 5$, nên $x + 3$ có thể bằng 3 hoặc 4.\n- Nếu $x + 3 = 3 \Rightarrow x = 0$.\n- Nếu $x + 3 = 4 \Rightarrow x = 1$.\nVậy có $2$ giá trị của $x$ thỏa mãn (đó là 0 và 1).
\textbf{Đáp số:} $2$.
}
\end{ex}

\begin{ex}
Ba số tự nhiên liên tiếp có tổng bằng 6. Hỏi số đứng ở giữa là số mấy?
\shortans{2}
\loigiai{
Ba số tự nhiên liên tiếp có tổng bằng 6 là: $1 + 2 + 3 = 6$.\nSố đứng ở giữa là số $2$.
\textbf{Đáp số:} $2$.
}
\end{ex}

\begin{ex}
Ba số tự nhiên liên tiếp có tổng bằng 9. Hỏi số lớn nhất trong ba số đó là số mấy?
\shortans{4}
\loigiai{
Ba số tự nhiên liên tiếp có tổng bằng 9 là: $2 + 3 + 4 = 9$.\nSố lớn nhất là số $4$.
\textbf{Đáp số:} $4$.
}
\end{ex}
